\section{How to View the Demonstration}
\label{sec:demo}

The demonstration can be viewed in three different ways.
The first way uses the hosted browser version and needs no installation.
The second way uses the desktop application, and the third way uses the command line tool.
Both of these need the build described in Section~\ref{sec:build}.

\subsection{Method 1. Browser Application}
\label{sec:browser}

The browser version runs the complete simulator inside the web browser, and it needs no server or installation.

\begin{enumerate}
\item Open \url{https://msrishab007.github.io/DigitalCircuitSimulator/} in a web browser.
\item Wait until the status bar reports that the engine is ready. The editor opens with an example circuit, as Figure~\ref{fig:browserhome} shows.
\item Choose an example from the \textbf{Load example} list, or type a netlist in the editor.
\item Select \textbf{Run simulation}. The Telemetry log tab reports the events processed, the critical path, and the result of each assertion.
\item Open the \textbf{Waveform viewer} tab to see the signals over time, as Figure~\ref{fig:browserwave} shows.
\item Open the \textbf{Schematic}, \textbf{Analysis}, and \textbf{How to use} tabs to see the circuit diagram, the analyses, and the user guide.
\end{enumerate}

\begin{figure}[H]
\centering
\screenshot{browser-home.png}{\textwidth}
\caption{Browser application after it loads}
\label{fig:browserhome}
\end{figure}

\begin{figure}[H]
\centering
\screenshot{browser-waveform.png}{\textwidth}
\caption{Waveform viewer in the browser application}
\label{fig:browserwave}
\end{figure}

The browser version can also run from the local copy of the source code.
This requires the build in Section~\ref{sec:build}, because the build creates the JavaScript engine.
The following commands serve the project folder, after which the page opens at \url{http://localhost:8000/}.

\begin{lstlisting}
cd DigitalCircuitSimulator
python -m http.server 8000
\end{lstlisting}

\subsection{Building the Project}
\label{sec:build}

The desktop application and the command line tool run from compiled classes.
The following commands download the source code and build the project.

\begin{lstlisting}
git clone https://github.com/MSRishab007/DigitalCircuitSimulator.git
cd DigitalCircuitSimulator
mvn clean package
\end{lstlisting}

The build compiles the code, runs the automated tests, measures coverage, and creates the browser engine.
A successful build ends with the message \texttt{BUILD SUCCESS}, as Figure~\ref{fig:build} shows.
The submitted source code folder can replace the \texttt{git clone} step.

\begin{figure}[H]
\centering
\screenshot{terminal-build.png}{\textwidth}
\caption{End of a successful Maven build}
\label{fig:build}
\end{figure}

\subsection{Method 2. Desktop Application}
\label{sec:desktop}

The desktop application offers the same features as the browser version in a Java Swing window.
The following command starts it from the project folder.

\begin{lstlisting}
java -cp target/classes simulator.gui.SimulatorDashboard
\end{lstlisting}

\begin{enumerate}
\item Choose an example from the \textbf{Load example} list, as Figure~\ref{fig:desktophome} shows. The editor on the left shows its netlist.
\item Select \textbf{Run simulation}. The Telemetry log tab reports the run.
\item Open the \textbf{Waveform viewer} tab. Figure~\ref{fig:desktopwave} shows the waveform of the two-bit counter.
\item Open the \textbf{Schematic} tab. The time slider colors each wire by its value. The \textbf{Collapse macros} and \textbf{Show critical path} options change the drawing, as Figure~\ref{fig:desktopschematic} shows.
\item Open the \textbf{Analysis} tab and select \textbf{Generate truth table}, \textbf{Analyze timing}, or \textbf{Fault coverage}. Figure~\ref{fig:desktopanalysis} shows an example.
\item Open the \textbf{How to use} tab to read the guide to the netlist language, as Figure~\ref{fig:desktopguide} shows.
\end{enumerate}

\begin{figure}[H]
\centering
\screenshot{desktop-home.png}{\textwidth}
\caption{Desktop application with the example list}
\label{fig:desktophome}
\end{figure}

\begin{figure}[H]
\centering
\screenshot{desktop-waveform.png}{\textwidth}
\caption{Waveform viewer in the desktop application}
\label{fig:desktopwave}
\end{figure}

\begin{figure}[H]
\centering
\screenshot{desktop-schematic.png}{\textwidth}
\caption{Schematic tab with the critical path highlighted}
\label{fig:desktopschematic}
\end{figure}

\begin{figure}[H]
\centering
\screenshot{desktop-analysis.png}{\textwidth}
\caption{Analysis tab in the desktop application}
\label{fig:desktopanalysis}
\end{figure}

\begin{figure}[H]
\centering
\screenshot{desktop-guide.png}{\textwidth}
\caption{User guide in the desktop application}
\label{fig:desktopguide}
\end{figure}

\subsection{Method 3. Command Line Tool}
\label{sec:cli}

The command line tool runs each operation as a separate command from the project folder.
The following command runs the half adder sample and writes its waveform to a VCD file.

\begin{lstlisting}
java -cp target/classes simulator.Main simulate samples/half_adder.netlist half_adder.vcd
\end{lstlisting}

The tool prints the progress of the run and the truth table of the circuit.

\begin{lstlisting}
1. Loading netlist: samples/half_adder.netlist
2. Validating topology...
3. Initializing simulation clock at zero...
4. Executing discrete event loop...
5. Generating output waveform...
Simulation complete. Processed 25 events.
Waveform written to half_adder.vcd
Truth table:
A B | Carry Sum
0 0 | 0     0
0 1 | 0     1
1 0 | 0     1
1 1 | 1     0
\end{lstlisting}

Table~\ref{tab:commands} lists the other commands.
Each command takes the same \texttt{java -cp target/classes simulator.Main} prefix.

\begin{table}[H]
\centering
\small
\caption{Commands of the command line tool}
\label{tab:commands}
\begin{tabular}{>{\raggedright\arraybackslash}p{6.6cm}>{\raggedright\arraybackslash}p{8.4cm}}
\toprule
\textbf{Command} & \textbf{Result} \\
\midrule
\texttt{simulate <netlist> [output.vcd]} & Runs the circuit and writes the waveform \\
\texttt{waveform-image <netlist> <output.png>} & Saves the waveform as an image \\
\texttt{validate <netlist>} & Checks the netlist without running it \\
\texttt{analyze <netlist>} & Reports the critical path and the maximum clock frequency \\
\texttt{truth-table <netlist>} & Prints the truth table or state table \\
\texttt{equivalent <first> <second>} & Compares two circuits \\
\texttt{fault-coverage <netlist>} & Reports the stuck-at faults that the stimuli detect \\
\texttt{check} & Runs the built-in self-check \\
\texttt{guide} & Prints the user guide \\
\bottomrule
\end{tabular}
\end{table}

The \texttt{check} command runs 22 end-to-end checks of the main features, as Figure~\ref{fig:check} shows.
It offers the quickest way to confirm that a build works.

\begin{lstlisting}
java -cp target/classes simulator.Main check
\end{lstlisting}

\begin{figure}[H]
\centering
\screenshot{terminal-check.png}{\textwidth}
\caption{Output of the self-check command}
\label{fig:check}
\end{figure}
